\documentclass[dvipdfmx,disablejfam]{jarticle}
\usepackage[dvipdfmx]{graphicx}
\usepackage{url}
\usepackage{amsmath}
\usepackage{amssymb}
\usepackage{bm}
\usepackage{float}
\usepackage{xcolor}

% 体裁（必要なら調整）
\topmargin -0.4mm
\oddsidemargin -0.4mm
\evensidemargin -0.4mm
\headheight 0mm
\headsep 0mm
\footskip 10mm
\textheight 247mm
\textwidth 170mm
\setlength{\parindent}{1zw}
\pagestyle{empty}

\makeatletter
\def \fgcaption{\def \@captype{figure} \caption}
\def \tbcaption{\def \@captype{table} \caption}
\makeatother

\def\figurename{Fig.}
\def\tablename{Table}
\def\topfraction{.95}
\def\bottomfraction{.95}
\def\textfraction{.05}
\def\floatpagefraction{.8}
\def\dbltopfraction{.95}
\def\dbltextfraction{.05}
\def\dblfloatpagefraction{.8}

\begin{document}

%==============================
%  ヘッダー
%==============================
\begin{center}
\begin{tabular}{|c|c|c|} \hline
  \parbox{50mm}{\vspace{1mm} \centering{\Large\bf 2026/04/16}} &
  \parbox{55mm}{\vspace{1mm} \centering{\Large\bf M1\quad 市川 詠祥}} &
  \parbox{50mm}{\vspace{1mm} \centering{\Large\bf No.~2}} \\ [2mm] \hline

  \multicolumn{3}{|c|}{
    \parbox{155mm}{
      \centering\Large\bf \vspace{2mm}
       くさび形モデルにおける表面波の解析
    }
  } \\ [3mm] \hline\hline

  \multicolumn{3}{|c|}{
    \parbox{155mm}{
      \large\vspace{2mm}

      {\Large\bf 今月の目標}
      \begin{itemize}
        \item くさび形のきずにおける表面波の解析
      \end{itemize}

      {\Large\bf 進捗の要点}
      \begin{itemize}
        \item くさび形きずにおける表面波シミュレーション
        \item 縦方向応力（T1）・横方向応力（T3）の動きを見るプログラムの実装
      \end{itemize}

      {\Large\bf 今後の実施事項}
      \begin{itemize}
        \item くさび形きずにおける表面波の解析
        \item 探触子の中心位置などを変えてみる
      \end{itemize}

      {\Large\bf 投稿・発表予定}
      \begin{itemize}
        \item
      \end{itemize}
    }
  } \\ [3mm] \hline
\end{tabular}
\end{center}

\vspace{6mm}
{\Large\bf 今回の内容}
\vspace{3mm}

\subsection*{1. 新入生向け発表（4/13）}
月曜日に新入生向けの研究紹介発表を行った．
先週末はその準備に充てていたため，シミュレーション関連の作業は月曜日以降から開始しました．

\subsection*{2. 表面波FDTDシミュレーションプログラムの作成（4/14）}
くさび形きずにおける表面波を観測するためのプログラム\texttt{kusabi\_surface\_sim.py}を新規作成しました．
このプログラムでは，縦方向応力$T_1$および横方向応力$T_3$を複数の観測点で時系列記録し，時空間マップとして保存できる構成とした．
また，探触子配置と計測点配置を確認するための可視化スクリプト（\texttt{kusabi\_T3\_map\_points.py}等）も合わせて作成した．
Fig.~\ref{fig:meas_points}に，くさび形表面における$T_3$計測点の配置確認図を示す．

\begin{figure}[H]
  \centering
  \includegraphics[width=0.85\textwidth, bb=0 0 640 302]{../../project4/tmp_output/kusabi_T3_map_pitch125_depth20.png}
  \caption{くさび形表面における$T_3$計測点の配置（pitch = 1.25 mm,depth = 0.20 mm）}
  \label{fig:meas_points}
\end{figure}

\subsection*{3. 時空間マップ可視化プログラムの実装（4/15）}
矩形型についての論文のFig.~3.5に相当する図を再現するため，可視化プログラム\path{plot_kusabi_T1_T3_surface_map.py}を新規作成した．
このプログラムは，シミュレーション結果（.npzファイル）を読み込み,
$T_1$（z方向の固定断面における縦方向応力の$x$-time マップ）と
$T_3$（表面計測点における横方向応力の$z$-time マップ）を
1枚の図にまとめて出力する,
あわせて\texttt{kusabi\_surface\_sim.py}のリファクタリングを行い，コードを整理した．
Fig.~\ref{fig:T1T3_map}に，得られた時空間マップを示す．

\begin{figure}[H]
  \centering
  \includegraphics[width=0.85\textwidth, bb=0 0 791 567]{../../project4/tmp_output/kusabi_T1_T3_surface_map_pitch125_depth20.png}
  \caption{くさび形きずにおける$T_1$・$T_3$時空間マップ（pitch = 1.25 mm, depth = 0.20 mm）}
  \label{fig:T1T3_map}
\end{figure}

%===========================================

\end{document}
